% Azimut2026 -- manuscript skeleton, version 0.1 (October 2026)
% Target journal (provisional): Solar Energy.
% This is a neutral article class for drafting, NOT the journal submission template.
% Change \drafttrue to \draftfalse before producing a clean reading copy.
\documentclass[11pt,a4paper]{article}

\usepackage[T1]{fontenc}
\usepackage[utf8]{inputenc}
\usepackage[english]{babel}
\usepackage{lmodern}
\usepackage{geometry}
\geometry{margin=2.6cm}
\usepackage{microtype}
\usepackage{amsmath,amssymb}
\usepackage{graphicx}
\usepackage{booktabs}
\usepackage{array}
\usepackage{enumitem}
\usepackage{xcolor}
\usepackage[numbers,sort&compress]{natbib}
\usepackage[hidelinks]{hyperref}
\usepackage{siunitx}

\newif\ifdraft
\drafttrue  % EDIT THIS: \draftfalse hides internal notes/checklist.
\newcommand{\draftnote}[1]{%
  \ifdraft\par\smallskip\noindent{\color{red!65!black}\sffamily\footnotesize[TO DO: #1]}\par\smallskip\fi
}
\newcommand{\figplaceholder}[1]{%
  \ifdraft\fbox{\begin{minipage}[c][3.4cm][c]{0.90\linewidth}
    \centering\sffamily\small Figure pending\\[2mm]#1
  \end{minipage}}%
  \else\fbox{\begin{minipage}[c][2cm][c]{0.90\linewidth}\centering Figure pending\end{minipage}}\fi
}
\newcommand{\GTI}{\mathrm{GTI}}
\newcommand{\GHI}{\mathrm{GHI}}
\newcommand{\DHI}{\mathrm{DHI}}
\newcommand{\DNI}{\mathrm{DNI}}
\newcommand{\RMSE}{\mathrm{RMSE}}
\newcommand{\MAE}{\mathrm{MAE}}
\newcommand{\MBE}{\mathrm{MBE}}

\title{Robust estimation of surface azimuth from tilted irradiance measurements:\\
experimental validation, uncertainty assessment, and field applications}
\author{Inti Piccioli\\[1mm]\small Coauthors and affiliations to be confirmed}
\date{Working draft -- October 2026}

\begin{document}
\maketitle

\begin{abstract}
An accurate description of the orientation of tilted irradiance sensors is essential
for the interpretation of global tilted irradiance (GTI) observations and for the
validation of irradiance transposition models. An approach based on the solar
hour angle of the daily GTI maximum was previously proposed to infer surface
azimuth from clear-sky observations. Here, we investigate the reliability of this
approach using a controlled measurement campaign with documented surface
orientations, and develop a framework for assessing fit quality, estimation
uncertainty, and robustness. The method is also examined as a tool for diagnosing
orientation inconsistencies in field measurements and quantifying their impact
on transposition-model evaluation.
\draftnote{Rewrite abstract when final numerical results exist: include sample size, azimuth range, main accuracy and uncertainty metrics, operational benefit, and limitations. Do not imply that field misalignments are independently verified unless they are.}
\end{abstract}

\noindent\textbf{Keywords:} tilted irradiance; surface azimuth; orientation
misalignment; pyranometer; uncertainty; quality control; transposition models

\section{Introduction}\label{sec:introduction}

\subsection{The importance of reliable surface-orientation metadata}

In situ measurements of global tilted irradiance (GTI) are essential to estimate the output of photovoltaic systems, and its accuracy is directly related to the economic risk of such project. Orientation errors in tilted sensors can propagate to irradiance transposition and model validation studies, particularly at unattended sites.

The optimal orientation of PV arrrays in terms of the year-round production is that of the Equator, i.e. due North (South) if the observer is located in the southern (northern) hemisphere. Small misalignments from this orientation can result in uncertainties that may impact in financial risk for PV projects. 

Misalignment of tilted sensors is something that can occur for a variety of reasons: incorrect determination of the true north-south direction, poor fixation of the sensors, extreme weather conditions that may move the platforms where the sensors are located.

Transposition models, that is, determination of GTI from GHI and assumptions about the diffuse component of radiation is highly dependent of the incidence angle of the direct component  of radiation, which also depends on the surface orientation, given by the tilt and azimuth.

\citep{menyhart2015,urraca2017}.
\draftnote{Add a quantitative motivation: how an azimuth error affects modeled GTI under representative conditions. Cite a dedicated study if used.}

\subsection{Existing approaches to orientation inference}
Orientation inference has been investigated using measured irradiance and
modeled clear-sky responses \citep{barbier2019,piccioli2021}, and using
photovoltaic generation measurements \citep{meng2020}. Our earlier conference
paper introduced the peak-hour-angle approach (PA) and compared it with a
model-error-minimization approach \citep{piccioli2021}. A subsequent thesis
applied orientation diagnosis in the context of transposition-model evaluation
\citep{piccioli2022thesis}.
\draftnote{Complete state-of-the-art search, including studies after 2021. Distinguish irradiance-sensor orientation from PV-system orientation; avoid claiming novelty for the original PA method.}

\subsection{Research gap and objectives}
The remaining challenge is not only to estimate an azimuth, but to determine
when that estimate is trustworthy. Specifically, the relationship between
measurement quality, stability of the fitted GTI maximum, uncertainty in
estimated azimuth, and practical usefulness for unattended stations requires
systematic assessment.

The objectives of this study are to (i) validate the peak-hour-angle approach
against documented surface orientations; (ii) quantify uncertainty and the
reliability of daily GTI fits; (iii) investigate robustness and multi-day
aggregation; and (iv) illustrate its potential for diagnosing orientation
inconsistencies and assessing their impact on transposition-model validation.
\draftnote{Tighten the contribution statement after analysis. Specify what is genuinely new compared with SWC2021 and the thesis.}

\section{Methodology}\label{sec:methodology}

\subsection{Solar geometry and azimuth convention}\label{subsec:geometry}
Let $\omega$ denote the solar hour angle and $\beta$ the surface tilt. We
define $\gamma=0$ for an equator-facing surface, with the sign convention
made explicit for the hemisphere under study. The daily hour angle of maximum
clear-sky GTI is denoted by $\omega^{*}$.
\draftnote{Add solar-angle equations, the exact implementation of solar time and timestamps, timezone convention, and the explicit signed azimuth definition (east/west) used by Azimut2026. Note how the convention maps to pvlib if applicable.}

\begin{figure}[t]
\centering
\figplaceholder{Geometry: tilt $\beta$, azimuth $\gamma$, hour angle $\omega$, and the controlled mount.}
\caption{Conceptual geometry and experimental configuration. Replace the
placeholder with a diagram and/or a photograph of the measurement system.}
\label{fig:geometry}
\end{figure}

\subsection{Mapping the hour angle of maximum GTI to surface azimuth}
\label{subsec:mapping}
For a selected day $d$, the modeled relationship between surface azimuth and
the location of the GTI maximum is represented by
\begin{equation}
  \hat\gamma_d = \mathcal{F}_d\!\left(\hat\omega_d^{*};\beta,\phi,\boldsymbol{p}_d\right),
  \label{eq:azimuth_mapping}
\end{equation}
where $\phi$ is the site latitude and $\boldsymbol{p}_d$ denotes the
assumptions and inputs used to construct the clear-sky/transposition mapping.
The 2021 formulation used a fitted mapping between $\omega^{*}$ and $\gamma$
\citep{piccioli2021}.
\draftnote{Describe the mapping ACTUALLY implemented in the repository (analytical fit, LUT, interpolator, clear-sky inputs, transposition model). Do not assume that the 2021 tangent fit is still used. Specify the azimuth search domain and numerical resolution.}

\subsection{Estimation of the daily GTI maximum}\label{subsec:dailyfit}
For a suitable interval around the daily maximum, a quadratic approximation
can be written as
\begin{equation}
  \GTI(\omega)\simeq a_2\omega^2+a_1\omega+a_0,
  \qquad
  \hat\omega^{*}=-\frac{a_1}{2a_2},\qquad a_2<0.
  \label{eq:quadratic}
\end{equation}
The location of the fitted maximum, rather than the largest individual
measurement, is used to reduce sensitivity to high-frequency fluctuations.
\draftnote{Specify fit-window definition, sample cadence, fitting algorithm, iteration/stopping rules, physical constraints, and treatment of non-concave or poorly constrained fits. Verify these against current code.}

\subsection{Clear-sky selection, outliers, and diagnostic indicators}
\label{subsec:quality}
A daily estimate is only meaningful if the irradiance maximum is sampled
adequately and the retained GTI observations are suitable for fitting.
Candidate diagnostics include fit residual dispersion, fraction of rejected
points, residual asymmetry, maximum curvature, and temporal data coverage.
\draftnote{Define the actual clear-sky and outlier masks with precise thresholds. Test the exceptional dates 2026-07-01, 2026-08-31, and 2026-09-17; distinguish physical cloud variability from an algorithmic selection artifact. Do not tune screening criteria solely to minimize observed angular error.}

\subsection{Uncertainty propagation and confidence assessment}
\label{subsec:uncertainty}
Writing $\mathbf{a}=(a_2,a_1,a_0)$, the formal local uncertainty of
$\hat\omega^{*}$ may be approximated by the delta method,
\begin{equation}
  u^2_{\omega^{*}}\simeq
  \left(\nabla_{\mathbf{a}}\omega^{*}\right)^{\mathsf{T}}
  \boldsymbol{\Sigma}_{\mathbf{a}}
  \nabla_{\mathbf{a}}\omega^{*},
  \label{eq:omega_uncertainty}
\end{equation}
with $\boldsymbol{\Sigma}_{\mathbf{a}}$ the covariance matrix of the
fitted coefficients. A first-order propagated component is
\begin{equation}
 u_{\gamma,\mathrm{fit}} \simeq
 \left|\frac{\partial\mathcal{F}_d}{\partial\omega^{*}}\right|
 u_{\omega^{*}}.
 \label{eq:gamma_uncertainty}
\end{equation}
This formal uncertainty does not, by itself, account for correlated residuals,
incorrect atmospheric inputs, timestamp offsets, or errors in the reference
orientation. Alternative uncertainty estimates and their calibration should
therefore be evaluated.
\draftnote{Compare covariance propagation with a time-aware bootstrap (e.g., blocks) and with empirical between-day spread. Report uncertainty components separately.}

\subsection{Multi-day estimation and temporal stability}
\label{subsec:aggregation}
For a period expected to have stable orientation, daily estimates
$\hat\gamma_d$ can be combined using robust or uncertainty-weighted procedures.
A weighted mean may be considered as a baseline,
\begin{equation}
  \hat\gamma_{\mathrm{period}}=
  \frac{\sum_d w_d\hat\gamma_d}{\sum_d w_d},
  \qquad w_d=u_{\gamma,d}^{-2},
  \label{eq:aggregation}
\end{equation}
provided that the daily uncertainties are calibrated and the effects of
correlations and shared systematic errors are addressed.
\draftnote{Compare median, robust location estimator, and weighted mean. Specify change-point detection procedure for historical stations, if used.}

\subsection{Validation protocol and evaluation metrics}
\label{subsec:metrics}
The azimuth error for controlled orientation $j$ is
$\Delta\gamma_j=\hat\gamma_j-\gamma_{\mathrm{ref},j}$.
Performance is described by bias, mean absolute error, and root mean square
error, together with repeatability across days:
\begin{equation}
  \MBE_{\gamma}=\frac{1}{N}\sum_{j=1}^{N}\Delta\gamma_j,\quad
  \MAE_{\gamma}=\frac{1}{N}\sum_{j=1}^{N}|\Delta\gamma_j|,\quad
  \RMSE_{\gamma}=\sqrt{\frac{1}{N}\sum_{j=1}^{N}\Delta\gamma_j^{2}}.
  \label{eq:metrics}
\end{equation}
\draftnote{Define unit of analysis (orientation period vs day), between-day dependence, confidence intervals, and whether reference orientation uncertainty is included. Establish an untouched test set or cross-period protocol before screening/model selection.}

\section{Data and experimental setup}\label{sec:data}

\subsection{Controlled azimuth campaign (2026)}\label{subsec:campaign}
The primary validation dataset consists of measurements of GTI acquired at multiple documented orientations of a fixed tilted surface. In order to achieve this, an inclined Kipp and Zonen CMP6 pyranometer, with a tilt of 35 degrees, was mounted on a rotating graduated plate located at the rooftop of the Faculty of Engineering in Montevideo, Uruguay (Latitude: $ \SI{-34.92}{\degree}$ , Longitude: $ \SI{-56.17}{\degree}$, Elevation: ???).  Special care was taken when leveling the table where the plate is mounted as well as the inclined plane where the pyranometer is placed.

GHI simultaneous data is acquired with a horizontal Kipp and Zonen (MODEL?) pyranometer, both calibrated and with traceability to the WMO (DETALLES?) and regularly maintained. The data is collected at 1-min intervals.

True north ($\gamma = \SI{0}{\degree}$) was obtained with the help of the shadow of a vertical rod under clear sky conditions at the solar noon, as provided by the NASA solar calculator (WEBSITE?). 
By comparing the north oriented GTI data under clear sky conditions with simultaneous GHI data, it was established that the misalignment of the tilted sensor was minimal  \draftnote{Assess uncertainty from the shadow thickness in cm, and discuss the gamma =0 correct orientation}


\begin{table}[t]
\centering
\caption{Provisional inventory of the 2026 controlled experiment. Fill in from
measurement metadata; do not infer reference precision from the step size.}
\label{tab:campaign}
\begin{tabular}{llll}
\toprule
Period / date & $\beta_{\mathrm{ref}}$ ($^\circ$) & $\gamma_{\mathrm{ref}}$ ($^\circ$) & Valid days \\
\midrule
To be completed & -- & -- & -- \\
\bottomrule
\end{tabular}
\end{table}

\subsection{Reference orientation and its uncertainty}
\label{subsec:reference}

The geographical North–South reference was established independently of the irradiance measurements using the shadow of a vertical rod at the predicted solar noon. The procedure was performed in late May using a cylindrical rod approximately 1 m high, mounted on a leveled horizontal surface. Rod verticality was checked against its supporting base using a right-angle reference.
The resulting meridian line was marked directly on the supporting table and used to establish the zero-azimuth reference of the rotating platform. The platform featured angular graduations at 1° intervals, a bubble level, and mechanical fixation screws.
The uncertainty associated with this reference procedure was explicitly considered when assessing absolute azimuth errors. Additionally, relative angular displacements were evaluated independently of the common initial alignment error.


We distinguish between nominal azimuth ($\gamma_{\mathrm{nom}}$), corresponding to the reported or intended surface orientation, and reference azimuth ($\gamma_{\mathrm{ref}}$), established through an independent geometric procedure. In the controlled experimental campaign, reference orientations were determined using an astronomically established North–South direction and a graduated rotating platform. For the remote stations, only nominal orientations were generally available, and the estimated azimuths were therefore interpreted as potential orientation corrections rather than independently validated measurements.
\draftnote{Document the reference procedure (survey, calibrated instrument, alignment, fixture), repeatability and estimated uncertainty for tilt and azimuth. Avoid claiming absolute angular accuracy better than the reference allows.}

Esto es un tema a discutir - Agus?


\subsection{Quality control and dataset construction}
\label{subsec:datasetqc}
All quality-control procedures should be specified before evaluating their
impact on azimuth errors. In addition to physical plausibility screening, the
availability of suitable clear-sky intervals around the GTI maximum must be
reported together with the fraction of excluded samples.
\draftnote{Create an auditable QC log: raw count -> daytime -> clear sky -> fit window -> outlier rejection -> valid estimates, with per-day reasons. State whether masks use GTI alone or additional channels.}


\subsection{Independent field applications}\label{subsec:fielddata}
To assess operational usefulness, the method will also be applied to
historical GTI measurements in which the reported orientation may be
uncertain or time-dependent \citep{piccioli2022thesis}.

\begin{itemize}
	\item Estación TA, describir equipos, timespan, medidas disponibles, instrumentación y describir evidencia de que la misma se encuentra desviada un ángulo indeterminado
	\item Esto resulta en modelos de transposición cuya validación es errónea, con métricas que no coinciden con lo esperado según la literatura
	\item GTI en cielo claro muestra un máximo en $\omega^{\ast} \neq 0$
	\item Reportar cantidad de días claros evaluados, valores obtenidos, promedio de azimuts
	\item Mostrar métricas con el valor de azimut calculado y ver cómo los porcentajes se acercan a lo esperado y conocido en la literatura.
\end{itemize}

\draftnote{Select and describe stations from the thesis; add site identifiers, available periods, instrumentation, and any independent evidence of orientation changes. State clearly which cases are diagnoses rather than externally validated ground truth.}

\section{Results}\label{sec:results}

\subsection{Dataset audit and behavior of the GTI fits}
\label{subsec:fitresults}
\draftnote{Show representative good and difficult fit examples (especially 2026-07-01, 2026-08-31, 2026-09-17). Quantify why observations are removed and how the rejection moves $\hat\omega^{*}$. Distinguish noisy data from biased fitting choices.}
\begin{figure}[t]
\centering
\figplaceholder{Panel examples: well-behaved curve; 2026-07-01; 2026-08-31; 2026-09-17.}
\caption{Representative fits to daily GTI maxima, including cases for which
measurement noise or sample rejection complicates the determination of
$\omega^{*}$.}
\label{fig:fits}
\end{figure}

\subsection{Accuracy and repeatability against controlled orientations}
\label{subsec:accuracy}
\draftnote{Provide overall $\MBE_\gamma$, $\MAE_\gamma$, $\RMSE_\gamma$, orientation-level uncertainty, and independent reference uncertainty. Plot $\hat\gamma$ versus $\gamma_{\rm ref}$, not only nominal values. Avoid pseudo-replication across days.}
\begin{figure}[t]
\centering
\figplaceholder{Estimated versus reference azimuth; 1:1 line; uncertainties; per-orientation markers.}
\caption{Estimated azimuth versus independently determined reference azimuth
for the controlled experiment.}
\label{fig:validation}
\end{figure}

\subsection{Fit quality, angular error, and uncertainty calibration}
\label{subsec:uncertaintyresults}
\draftnote{Evaluate correlation/predictive value of RMSE, reject fraction, asymmetry, curvature, and $u_{\omega^{*}}$; compare predicted intervals with observed errors. Do not describe a proxy as calibrated uncertainty without coverage testing.}
\begin{figure}[t]
\centering
\figplaceholder{Angular error and repeatability versus azimuth; fitted uncertainty and empirical coverage.}
\caption{Error distribution and reliability of angular uncertainty estimates
as a function of surface orientation and GTI-fit diagnostics.}
\label{fig:uncertainty}
\end{figure}

\subsection{Benefits of multi-day aggregation}
\label{subsec:multidayresults}
\draftnote{Show convergence of error and between-sample variability as a function of the number of independent clear-sky days. Compare robust estimators and predefine resampling.}
\begin{figure}[t]
\centering
\figplaceholder{Error/dispersion versus number of clear days and selected perturbation tests.}
\caption{Convergence of multi-day azimuth estimation and sensitivity to data
availability or fitting conditions.}
\label{fig:robustness}
\end{figure}

\subsection{Sensitivity to measurement and modeling perturbations}
\label{subsec:sensitivity}
\draftnote{Perturb tilt, time stamps, cadence, GTI noise, missing data, fit interval and clear-sky inputs. Prioritize confounding between apparent azimuth error and timestamp/tilt errors. Separate simulation from field evidence. Supplemental plots as needed.}

\subsection{Diagnosis of field-station orientation}
\label{subsec:fieldresults}
\draftnote{Show $\hat\gamma(t)$ and uncertainty by station, stable periods and potential changes, with any external supporting evidence. Avoid asserting true azimuth from GTI alone.}
\begin{figure}[t]
\centering
\figplaceholder{Time series of field-station estimated azimuth and documented station events.}
\caption{Application to historical field-station measurements: temporal
behavior and possible orientation changes.}
\label{fig:field}
\end{figure}

\subsection{Consequences for transposition-model evaluation}
\label{subsec:transpositionresults}
\draftnote{Compare $\GTI_{\rm model}(\gamma_{\rm nom})$ and $\GTI_{\rm model}(\hat\gamma)$ with measurements. Report MBD/RMSD and number of samples, on held-out data if possible. Test multiple models and discuss non-independence if the mapping itself uses a transposition model.}
\begin{table}[t]
\centering
\caption{Illustrative layout for the impact of orientation metadata on
transposition-model evaluation. All entries remain to be computed.}
\label{tab:transposition}
\begin{tabular}{lllll}
\toprule
Station & Model & Azimuth source & MBD & RMSD \\
\midrule
-- & -- & Nominal & -- & -- \\
-- & -- & Estimated & -- & -- \\
\bottomrule
\end{tabular}
\end{table}

\section{Discussion}\label{sec:discussion}

\subsection{When is a daily azimuth estimate trustworthy?}
\draftnote{Summarize practical acceptance criteria, uncertainty calibration, and minimum clear-day requirements. Separate recommendations supported by controlled data from extrapolations.}

\subsection{Identifiability and potential confounding factors}
\draftnote{Address timestamp shifts, incorrect tilt, horizon obstructions, partial shading, directional/cosine response, soiling, atmospheric variability, low curvature, model dependence, and seasonal effects. Clarify the meaning of `effective azimuth'.}

\subsection{Relation to earlier work and scope of the contribution}
The present work builds on the previously published peak-hour-angle method
\citep{piccioli2021} and its field application in a thesis
\citep{piccioli2022thesis}. Its intended contribution is a more rigorous
experimental validation and an explicit reliability assessment, rather than
re-introduction of the original geometric method.
\draftnote{Compare fairly with Barbier et al. and Meng et al. State differences in required inputs, data type, scope, and dependence on transposition models. Make clear what results are genuinely new.}

\subsection{Implications for solar-resource assessment}
\draftnote{Explain when orientation diagnosis can improve metadata QC and model assessment. Avoid conflating reduced model error with proof that the physical azimuth is exact.}

\section{Conclusions}\label{sec:conclusions}
\draftnote{Final conclusions should contain 3-4 quantitative claims: angular performance on the controlled dataset; calibrated reliability criteria; robustness limits; demonstrated consequences for transposition models. Include clear limitations (tilt range, reference uncertainty, site/season coverage).}
This section will be completed after consolidation of the controlled
validation and field-application results.


\section*{Data and code availability}
\draftnote{Specify public repository/version, persistent identifier if available, access to quality-controlled data, and limitations on redistributing third-party datasets.}
To be completed.

\section*{Acknowledgements}
\draftnote{Funding, institutions, measurement technicians, collaborators, and computational infrastructure.}
To be completed.

\bibliographystyle{unsrtnat}
\bibliography{references}

\ifdraft
\clearpage
\appendix
\section{Internal research checklist (remove before submission)}
This appendix is visible only while \texttt{\\drafttrue} is enabled. The
more detailed, prioritized roadmap is in \texttt{TODO.md}.
\begin{enumerate}[leftmargin=*,label=\textbf{\arabic*.}]
\item \textbf{Experimental inventory:} metadata, reference orientation,
  timestamp control, sensor specifications, quality-control flow diagram.
\item \textbf{GTI fits:} audit difficult dates (2026-07-01, 2026-08-31,
  2026-09-17), quantify removed samples and shifts in $\hat\omega^{*}$.
\item \textbf{Uncertainty:} propagate fit covariance, assess time-aware
  bootstrap, test empirical coverage of angular intervals.
\item \textbf{Core validation:} error by orientation, repeatability, and
  independence of evaluation units.
\item \textbf{Multi-day estimation:} convergence and change-point sensitivity.
\item \textbf{Robustness:} tilt/timestamp offsets, noise, missing data,
  fitting window, temporal resolution, and modeling assumptions.
\item \textbf{Field application:} recover thesis stations and document
  available independent orientation evidence.
\item \textbf{Transposition:} nominal versus inferred azimuth using a
  holdout period and more than one model when feasible.
\item \textbf{Writing:} update literature; insert real figures and tables;
  complete abstract, limitations, reproducibility and acknowledgements.
\end{enumerate}

\fi
\end{document}
